\documentclass[11pt,letterpaper]{article}
\usepackage[utf8]{inputenc}
\usepackage[spanish]{babel}
\usepackage[margin=2.2cm]{geometry}
\usepackage{enumitem}
\usepackage{listings}
\usepackage{xcolor}
\usepackage{titlesec}
\usepackage{fancyhdr}
\usepackage{parskip}

% ── Colores ──────────────────────────────────────────────
\definecolor{bg}{HTML}{F8F9FA}
\definecolor{kw}{HTML}{2563EB}
\definecolor{cm}{HTML}{6B7280}
\definecolor{str}{HTML}{059669}
\definecolor{accent}{HTML}{D97706}

% ── Código Python ────────────────────────────────────────
\lstset{
  language=Python,
  backgroundcolor=\color{bg},
  basicstyle=\ttfamily\small,
  keywordstyle=\color{kw}\bfseries,
  commentstyle=\color{cm}\itshape,
  stringstyle=\color{str},
  showstringspaces=false,
  frame=single,
  rulecolor=\color{black!20},
  breaklines=true,
  xleftmargin=6pt,
  xrightmargin=6pt,
  aboveskip=8pt,
  belowskip=8pt,
  literate={á}{{\'a}}1 {é}{{\'e}}1 {í}{{\'i}}1 {ó}{{\'o}}1 {ú}{{\'u}}1 {ñ}{{\~n}}1
}

% ── Encabezado ───────────────────────────────────────────
\pagestyle{fancy}
\fancyhf{}
\rhead{\small Programación Avanzada}
\lhead{\small POO en Python --- Ejercicios}
\rfoot{\small Página \thepage}
\renewcommand{\headrulewidth}{0.4pt}

% ── Títulos ──────────────────────────────────────────────
\titleformat{\section}{\large\bfseries\color{accent}}{Ejercicio \thesection.}{6pt}{}
\titleformat{\subsection}[runin]{\bfseries}{}{0pt}{}[~~---]

% ════════════════════════════════════════════════════════════
\begin{document}

\begin{center}
  {\LARGE\bfseries Ejercicios de POO en Python}\\[6pt]
  {\large Clase 03 --- Clases, objetos, herencia y composición}\\[4pt]
  {\color{cm}\small 10 ejercicios $\cdot$ Dificultad: baja $\to$ media $\to$ alta}
\end{center}

\vspace{.6cm}

% ─────────────────────────────────────────────────────────
%  BAJA
% ─────────────────────────────────────────────────────────

\hfill\fbox{\small\bfseries DIFICULTAD BAJA}
\vspace{2pt}

% ── 1 ────────────────────────────────────────────────────
\section{Rectángulo}
\subsection{Objetivo} Crear una clase básica con atributos y métodos simples.

Crea la clase \texttt{Rectangulo} con:
\begin{itemize}[nosep]
  \item \textbf{Atributos:} \texttt{base: float}, \texttt{altura: float}
  \item \textbf{Métodos:}
  \begin{itemize}[nosep]
    \item \texttt{area() -> float} --- retorna base $\times$ altura
    \item \texttt{perimetro() -> float} --- retorna $2 \times (\text{base} + \text{altura})$
    \item \texttt{\_\_str\_\_} $\to$ \texttt{"Rectángulo 5.0 x 3.0"}
  \end{itemize}
\end{itemize}

\textbf{Prueba esperada:}
\begin{lstlisting}
r = Rectangulo(5, 3)
print(r.area())       # 15.0
print(r.perimetro())  # 16.0
print(r)              # Rectángulo 5.0 x 3.0
\end{lstlisting}

% ── 2 ────────────────────────────────────────────────────
\section{Contador}
\subsection{Objetivo} Practicar métodos que modifican estado interno.

Crea la clase \texttt{Contador} con:
\begin{itemize}[nosep]
  \item \textbf{Atributos:} \texttt{\_valor: int} (inicia en 0)
  \item \textbf{Métodos:}
  \begin{itemize}[nosep]
    \item \texttt{incrementar()} --- suma 1 al valor
    \item \texttt{decrementar()} --- resta 1, pero nunca baja de 0
    \item \texttt{get\_valor() -> int}
    \item \texttt{resetear()} --- vuelve el valor a 0
  \end{itemize}
\end{itemize}

\textbf{Prueba esperada:}
\begin{lstlisting}
c = Contador()
c.incrementar()
c.incrementar()
c.incrementar()
print(c.get_valor())  # 3
c.decrementar()
print(c.get_valor())  # 2
c.resetear()
print(c.get_valor())  # 0
\end{lstlisting}

% ── 3 ────────────────────────────────────────────────────
\section{Mascota}
\subsection{Objetivo} Usar \texttt{\_\_init\_\_} con valores por defecto y \texttt{\_\_str\_\_}.

Crea la clase \texttt{Mascota} con:
\begin{itemize}[nosep]
  \item \textbf{Atributos:} \texttt{nombre: str}, \texttt{especie: str}, \texttt{edad: int} (por defecto 0)
  \item \textbf{Métodos:}
  \begin{itemize}[nosep]
    \item \texttt{cumpleanios()} --- suma 1 a la edad
    \item \texttt{es\_adulta() -> bool} --- \texttt{True} si edad $\geq$ 3
    \item \texttt{\_\_str\_\_} $\to$ \texttt{"Luna (Gato) - 2 años"}
  \end{itemize}
\end{itemize}

\textbf{Prueba esperada:}
\begin{lstlisting}
m = Mascota("Luna", "Gato")
m.cumpleanios()
m.cumpleanios()
print(m)              # Luna (Gato) - 2 años
print(m.es_adulta())  # False
m.cumpleanios()
print(m.es_adulta())  # True
\end{lstlisting}

% ─────────────────────────────────────────────────────────
%  MEDIA
% ─────────────────────────────────────────────────────────

\vspace{.5cm}
\hfill\fbox{\small\bfseries DIFICULTAD MEDIA}
\vspace{2pt}

% ── 4 ────────────────────────────────────────────────────
\section{Inventario de Tienda}
\subsection{Objetivo} Trabajar con listas de objetos dentro de otro objeto.

Crea las clases \texttt{Producto} e \texttt{Inventario}:

\texttt{Producto}:
\begin{itemize}[nosep]
  \item \textbf{Atributos:} \texttt{nombre: str}, \texttt{precio: float}, \texttt{cantidad: int}
  \item \textbf{Método:} \texttt{valor\_total() -> float} --- retorna precio $\times$ cantidad
\end{itemize}

\texttt{Inventario}:
\begin{itemize}[nosep]
  \item \textbf{Atributos:} \texttt{\_productos: list} (vacía al inicio)
  \item \textbf{Métodos:}
  \begin{itemize}[nosep]
    \item \texttt{agregar(producto: Producto)}
    \item \texttt{valor\_inventario() -> float} --- suma el \texttt{valor\_total()} de todos los productos
    \item \texttt{producto\_mas\_caro() -> str} --- retorna el nombre del producto con mayor precio unitario
  \end{itemize}
\end{itemize}

\textbf{Prueba esperada:}
\begin{lstlisting}
inv = Inventario()
inv.agregar(Producto("Mouse", 25.0, 10))
inv.agregar(Producto("Teclado", 45.0, 5))
inv.agregar(Producto("Monitor", 200.0, 3))
print(inv.valor_inventario())  # 1075.0
print(inv.producto_mas_caro()) # Monitor
\end{lstlisting}

% ── 5 ────────────────────────────────────────────────────
\section{Temperatura con @property}
\subsection{Objetivo} Usar \texttt{@property} y setter con validación.

Crea la clase \texttt{Temperatura} con:
\begin{itemize}[nosep]
  \item \textbf{Atributo privado:} \texttt{\_celsius: float}
  \item \textbf{Properties:}
  \begin{itemize}[nosep]
    \item \texttt{celsius} (getter y setter) --- el setter no permite valores menores a $-273.15$
    \item \texttt{fahrenheit} (getter, solo lectura) --- retorna $\text{celsius} \times 9/5 + 32$
  \end{itemize}
  \item \textbf{Método:} \texttt{\_\_str\_\_} $\to$ \texttt{"25.0°C / 77.0°F"}
\end{itemize}

\textbf{Prueba esperada:}
\begin{lstlisting}
t = Temperatura(25)
print(t.celsius)     # 25
print(t.fahrenheit)  # 77.0
print(t)             # 25.0°C / 77.0°F
t.celsius = -300     # No cambia (menor al cero absoluto)
print(t.celsius)     # 25
\end{lstlisting}

% ── 6 ────────────────────────────────────────────────────
\section{Figuras con herencia}
\subsection{Objetivo} Herencia, \texttt{super()} y polimorfismo con override.

Crea una clase base \texttt{Figura} y dos clases hijas:

\texttt{Figura}:
\begin{itemize}[nosep]
  \item \textbf{Atributo:} \texttt{color: str}
  \item \textbf{Método:} \texttt{area() -> float} --- retorna 0 (será sobreescrito)
  \item \textbf{Método:} \texttt{\_\_str\_\_} $\to$ \texttt{"Figura color rojo, área: 0"}
\end{itemize}

\texttt{Circulo(Figura)}:
\begin{itemize}[nosep]
  \item \textbf{Atributo adicional:} \texttt{radio: float}
  \item \textbf{Override:} \texttt{area()} $\to$ $\pi \times r^2$
\end{itemize}

\texttt{Triangulo(Figura)}:
\begin{itemize}[nosep]
  \item \textbf{Atributos adicionales:} \texttt{base: float}, \texttt{altura: float}
  \item \textbf{Override:} \texttt{area()} $\to$ $\text{base} \times \text{altura} / 2$
\end{itemize}

\textbf{Prueba esperada:}
\begin{lstlisting}
figuras = [Circulo("rojo", 5), Triangulo("azul", 10, 4)]
for f in figuras:
    print(f)
# Figura color rojo, área: 78.54
# Figura color azul, área: 20.0
\end{lstlisting}

% ── 7 ────────────────────────────────────────────────────
\section{Concatenación entre objetos}
\subsection{Objetivo} Usar el resultado de un método como argumento de otro objeto.

Crea las clases \texttt{Conversor} y \texttt{Formateador}:

\texttt{Conversor}:
\begin{itemize}[nosep]
  \item \texttt{km\_a\_millas(km: float) -> float} --- retorna $\text{km} \times 0.621371$
  \item \texttt{kg\_a\_libras(kg: float) -> float} --- retorna $\text{kg} \times 2.20462$
\end{itemize}

\texttt{Formateador}:
\begin{itemize}[nosep]
  \item \texttt{mostrar(valor: float, unidad: str) -> str} --- retorna \texttt{"15.53 mi"} (redondeado a 2 decimales)
\end{itemize}

\textbf{Prueba esperada (concatenado en una sola línea):}
\begin{lstlisting}
conv = Conversor()
fmt  = Formateador()

# Concatenado:
print(fmt.mostrar( conv.km_a_millas(25), "mi" ))
# 15.53 mi

print(fmt.mostrar( conv.kg_a_libras(70), "lb" ))
# 154.32 lb
\end{lstlisting}

% ─────────────────────────────────────────────────────────
%  ALTA
% ─────────────────────────────────────────────────────────

\vspace{.5cm}
\hfill\fbox{\small\bfseries DIFICULTAD ALTA}
\vspace{2pt}

% ── 8 ────────────────────────────────────────────────────
\section{Sistema de Empleados}
\subsection{Objetivo} Herencia múltiple con \texttt{super()}, override y composición.

Crea la siguiente jerarquía:

\texttt{Empleado}:
\begin{itemize}[nosep]
  \item \textbf{Atributos:} \texttt{nombre: str}, \texttt{\_sueldo\_base: float}
  \item \textbf{Métodos:}
  \begin{itemize}[nosep]
    \item \texttt{calcular\_sueldo() -> float} --- retorna el sueldo base
    \item \texttt{\_\_str\_\_} $\to$ \texttt{"Ana - Sueldo: \$3000.0"}
  \end{itemize}
\end{itemize}

\texttt{Gerente(Empleado)}:
\begin{itemize}[nosep]
  \item \textbf{Atributo adicional:} \texttt{bono: float}
  \item \textbf{Override:} \texttt{calcular\_sueldo()} $\to$ sueldo base + bono
\end{itemize}

\texttt{Vendedor(Empleado)}:
\begin{itemize}[nosep]
  \item \textbf{Atributos adicionales:} \texttt{ventas: float}, \texttt{comision: float} (porcentaje, ej: 0.05)
  \item \textbf{Override:} \texttt{calcular\_sueldo()} $\to$ sueldo base + (ventas $\times$ comisión)
\end{itemize}

Crea también la clase \texttt{Nomina}:
\begin{itemize}[nosep]
  \item \textbf{Atributos:} \texttt{\_empleados: list}
  \item \textbf{Métodos:}
  \begin{itemize}[nosep]
    \item \texttt{agregar(emp: Empleado)}
    \item \texttt{total\_nomina() -> float} --- suma \texttt{calcular\_sueldo()} de todos (polimorfismo)
    \item \texttt{reporte()} --- imprime cada empleado con su sueldo
  \end{itemize}
\end{itemize}

\textbf{Prueba esperada:}
\begin{lstlisting}
n = Nomina()
n.agregar(Empleado("Ana", 3000))
n.agregar(Gerente("Luis", 4000, 1500))
n.agregar(Vendedor("María", 2500, 50000, 0.05))

n.reporte()
# Ana - Sueldo: $3000.0
# Luis - Sueldo: $5500.0
# María - Sueldo: $5000.0

print(n.total_nomina())  # 13500.0
\end{lstlisting}

% ── 9 ────────────────────────────────────────────────────
\section{Banco con transferencias}
\subsection{Objetivo} Un objeto recibe otro objeto como argumento de un método.

Reutiliza la clase \texttt{CuentaBancaria} del ejercicio de clase (o créala de nuevo) y agrégale el método:
\begin{itemize}[nosep]
  \item \texttt{transferir(destino: CuentaBancaria, monto: float) -> bool}
  \begin{itemize}[nosep]
    \item Retira \texttt{monto} de la cuenta actual
    \item Si el retiro fue exitoso, deposita en \texttt{destino}
    \item Retorna \texttt{True} si se completó, \texttt{False} si no había fondos
  \end{itemize}
\end{itemize}

\textbf{Prueba esperada:}
\begin{lstlisting}
c1 = CuentaBancaria("Ana", 5000)
c2 = CuentaBancaria("Luis", 1000)

ok = c1.transferir(c2, 2000)
print(ok)               # True
print(c1.get_saldo())   # 3000
print(c2.get_saldo())   # 3000

fail = c1.transferir(c2, 99999)
print(fail)             # False
print(c1.get_saldo())   # 3000  (no cambió)
\end{lstlisting}

% ── 10 ───────────────────────────────────────────────────
\section{Cadena de procesamiento}
\subsection{Objetivo} Concatenar múltiples objetos de distintas clases en una sola expresión.

Crea tres clases que trabajen juntas:

\texttt{Sensor}:
\begin{itemize}[nosep]
  \item \textbf{Atributo:} \texttt{nombre: str}
  \item \textbf{Método:} \texttt{leer() -> float} --- retorna un valor fijo (simula una lectura, ej: 36.7)
\end{itemize}

\texttt{Procesador}:
\begin{itemize}[nosep]
  \item \textbf{Método:} \texttt{normalizar(valor: float, minimo: float, maximo: float) -> float}
  \begin{itemize}[nosep]
    \item Retorna $(\text{valor} - \text{mínimo}) / (\text{máximo} - \text{mínimo})$
  \end{itemize}
\end{itemize}

\texttt{Alerta}:
\begin{itemize}[nosep]
  \item \textbf{Método:} \texttt{evaluar(valor\_normalizado: float) -> str}
  \begin{itemize}[nosep]
    \item Si $< 0.3$: retorna \texttt{"BAJO"}
    \item Si $< 0.7$: retorna \texttt{"NORMAL"}
    \item Si $\geq 0.7$: retorna \texttt{"ALTO"}
  \end{itemize}
\end{itemize}

\textbf{Prueba esperada (todo concatenado):}
\begin{lstlisting}
s = Sensor("Temperatura")   # leer() retorna 36.7
p = Procesador()
a = Alerta()

# Concatenado en una sola línea:
estado = a.evaluar( p.normalizar( s.leer(), 35.0, 42.0 ) )
print(estado)  # NORMAL

# Equivalente paso a paso:
lectura     = s.leer()                         # 36.7
normalizado = p.normalizar(lectura, 35.0, 42.0) # 0.2428...
estado      = a.evaluar(normalizado)            # NORMAL
\end{lstlisting}

\vspace{1cm}
\begin{center}
  \rule{.4\textwidth}{.4pt}\\[6pt]
  {\small\color{cm} Todos los ejercicios deben incluir \texttt{type hints} en parámetros y retornos.}
\end{center}

\end{document}
